% 从旧数学附录迁出的莫队二次离线说明；不新增算法。
\section{莫队二次离线：前缀事件与自身扣除}\label{knowledge-mo-secondary}
对区间内满足 $H(i,j)=[\operatorname{popcount}(a_i\mathbin{\mathrm{xor}}a_j)=k]$
的无序下标对，使用 \texttt{xor\_hamming\_pairs}。
令 $t(i)=\sum_{j<i}H(i,j)$，$P(x)=\sum_{i<x}t(i)$，
以及 $T(p,i)=\sum_{j<p}H(i,j)-[k=0\land i<p]$。
右端扩展时贡献为 $t(i)-T(L,i)$，左端扩展为 $T(R,i)-t(i)$；
删除取相反数。整段端点移动的 $t$ 和由 $P$ 的区间差求出，
剩余部分存成按前缀长度分桶的事件。
第二遍顺扫只做插入，每次枚举异或后有 $k$ 个置位的掩码更新计数；
事件中的每个下标只读一次计数，并在 $k=0$ 时扣自身。
最后按莫队顺序累加窗口答案差，再按原编号返回。

区间统一为 0-based 半开区间，先扩右、扩左，再缩右、缩左，保持窗口合法。
允许空区间；没有查询时直接返回。位宽默认为 14，若 $k$ 超出位宽则答案全零。
P4887 的 $k<2^{14}$ 并不代表 $k\le14$。
令 $C=\binom{b}{k}$，块长约为 $n/\sqrt m$，
总时间 $O(2^b+nC+m\log(m+1)+n\sqrt m+m)$，空间 $O(2^b+n+m)$。
端点移动不再乘枚举掩码数 $C$。答案与前缀和使用 long long；
十万个相等数的配对数超过 int32。
本地全枚举、随机逐对参考及十万规模的 $k=0,7$ 分支通过；线上验证待补。

